\documentclass[11pt]{article}
\usepackage{amsmath,amssymb}
\usepackage{xurl}
\usepackage[hidelinks]{hyperref}
% Shared commands for notes in docs/paper-gaps/.
%
% Two halves.  The link commands below are project identity: OWNER/REPO is the
% placeholder that scripts/bootstrap_project.py replaces with this project's
% GitHub slug and Pages site.  Everything after them is NOTATION, and notation
% belongs to the source paper: the definitions here are an example set, and a
% project replaces them with the macros of its own paper's style file.  The one
% rule that never changes: a command defined here must mean exactly what the
% source paper's macro means, or a note silently says something else.

\newcommand{\Lean}{\textsc{Lean}}
\newcommand{\leanid}[1]{\textsf{#1}}
\newcommand{\ghissue}[1]{\href{https://github.com/OWNER/REPO/issues/#1}{GitHub issue \##1}}
\newcommand{\ghpr}[1]{\href{https://github.com/OWNER/REPO/pull/#1}{PR \##1}}
% Local issue tracker (issues/NNNN-*.md in this repository); no remote URL.
\newcommand{\localissue}[1]{issue \##1 (\path{issues/})}
\newcommand{\doclink}[1]{\href{https://OWNER.github.io/REPO/docs/#1.html}{\path{#1}}}

% --- Notation: replace with the source paper's own macros --------------------
% \F, \N, \C, \R, \tr, \ot, \eps, \E, \bx, \bu, \bv, \by

\newcommand{\F}{\mathbb{F}}
\newcommand{\N}{\mathbb{N}}
\newcommand{\C}{\mathbb{C}}
\newcommand{\R}{\mathbb{R}}
\newcommand{\tr}{\mathrm{tr}}
\newcommand{\eps}{\epsilon}
\newcommand{\ot}{\otimes}
\newcommand{\E}{\mathop{\bf E\/}}

% bold random variables
\newcommand{\bu}{\boldsymbol{u}}
\newcommand{\bv}{\boldsymbol{v}}
\newcommand{\bx}{{\boldsymbol{x}}}
\newcommand{\by}{\boldsymbol{y}}

% T1 so literal < and > in \texttt placeholders typeset as themselves
\usepackage[T1]{fontenc}

% bra-ket
\usepackage{braket}

% --- Example structured spaces ---
\newcommand{\polyfunc}[3]{\mathcal{P}(#1, #2, #3)}
\newcommand{\polymeas}[3]{\mathrm{PolyMeas}(#1, #2, #3)}
\newcommand{\polysub}[3]{\mathrm{PolySub}(#1, #2, #3)}

% Approximation relations (paper uses \approx_\eps and \simeq_\zeta)
\newcommand{\approxeps}{\approx_{\varepsilon}}
\newcommand{\simgeq}{\simeq_{\zeta}}


\title{Index Pairing in the Low-Degree Sandwich}
\author{}
\date{2026-09-03}

\begin{document}
\maketitle

\section{At a glance}
\begin{itemize}
  \item \textbf{Difficulty.} Low. The source has a typographical mismatch
  between a measurement family and the outcome supplied to it.
  \item \textbf{Estimated weight.} The correction changes only the displayed
  index pairing; the imported proof applies without a new estimate.
  \item \textbf{Mathlib/project split.} No Mathlib result is needed to identify
  the mismatch. The formal statement uses a dependent family of outcome types.
  \item \textbf{Key mathematical inputs.} Each projector must receive an
  outcome from its own alphabet, and an empty ordered product equals the
  identity. The correction is tracked in
  \href{https://github.com/Dengnifer/MIPStarRE-A/issues/16}{GitHub issue \#16}
  and descends from \localissue{0002}.
\end{itemize}

\section{Key theorem forms: the printed statement}
For $i\in\{1,\ldots,k\}$, let $G^{i,x}$ be a projective measurement with
outcome set $\mathcal G_i$. The source prints the sandwich as
\[
  G^{k,x}_{g_1}\cdots G^{1,x}_{g_k}\cdots G^{k,x}_{g_1}.
\]
This expression is not defined for a general family of outcome sets: the
outer factor belongs to the measurement indexed by $\mathcal G_k$, whereas
$g_1$ belongs to $\mathcal G_1$. The mismatch occurs in
\cite[lines~465--501, especially lines~484--488]{Ji2020MIPStar}.

The result is imported there from Fact~4.34 of the NEEXP paper
\cite[lines~952--973, especially lines~965--968]{Natarajan2018NEEXP}, mirrored
at \path{references/neexp-paper/05_quantum_preliminaries.tex}. That statement
pairs every measurement with its own outcome, giving
\[
  G^{k,x}_{g_k}\cdots G^{2,x}_{g_2}G^{1,x}_{g_1}
  G^{2,x}_{g_2}\cdots G^{k,x}_{g_k}.
\]
The order of the operators is unchanged; only the subscripts are corrected.

\section{Blueprint and formalization}
The blueprint uses the type-correct pairing at \path{lem:ld-sandwich} in
\path{blueprint/src/chapter/ch12_qpbt_games.tex} and explains the discrepancy
in \path{rem:ld-sandwich-indexing}. The formal definition
\leanid{sandwichProduct} uses an outcome function whose value at $i$ belongs
to the outcome type of $G^i$. Thus every operator receives the corresponding
component $g_i$. When $k=0$, the product is empty and its value is the identity,
which supplies the unique empty-tuple outcome.

\section{Conclusion}
The reversed outcome subscripts in the source are a typographical error, not a
different mathematical construction. The blueprint and Lean statement retain
the source's operator order while restoring the only type-correct pairing.

\bibliographystyle{alpha}
\bibliography{references}
\end{document}
